% Complete unchanged phase258 programs; no omitted diagnostics or controls.
\documentclass[11pt,a4paper]{article}
\usepackage[margin=1in]{geometry}
\usepackage{fontspec}
\setmainfont{STIXGeneral}
\setmonofont{Menlo}
\usepackage{fvextra}
\usepackage{hyperref}
\DefineVerbatimEnvironment{CompleteProgram}{Verbatim}{fontsize=\small,breaklines,breakanywhere,breakanywheresymbolpre={},breakanywheresymbolpost={},breaksymbolleft={},breaksymbolright={}}
\setlength{\parindent}{0pt}
\begin{document}
\section*{Complete additional programs for projective-line maniplexes}
The following two listings contain the complete unchanged programs used by
the admitted application. The shared projective-line finite program is charged
once in its separate appendix. No control or diagnostic line has been removed.
\subsection*{Local string markings}
\noindent\texttt{phase258\_local\_sggi.g}
% BEGIN COMPLETE FILE phase258_local_sggi.g
\begin{CompleteProgram}
# All string generating triples, without assuming intersection property.
Need258:=function(b,s) if not b then Error(s);fi;end;
IP258:=function(G,rr)
 local ss,hs,i,j;
 ss:=Combinations([1..Length(rr)]);hs:=List(ss,s->Subgroup(G,rr{s}));
 for i in [1..Length(ss)] do for j in [1..Length(ss)] do
 if Intersection(hs[i],hs[j])<>hs[Position(ss,Intersection(ss[i],ss[j]))] then return false;fi;
 od;od;return true;
end;
Run258:=function()
 local pair,G,inv,a,b,c,rr,total,bad,types,t,markings,aut,left,orb,cs,H,z,x,y;
 for pair in [[SymmetricGroup(4),72,[[3,3,4],[3,4,3],[4,3,3]]],
              [AlternatingGroup(5),360,[[3,5,5],[5,3,5],[5,5,3]]]] do
 G:=pair[1];inv:=Filtered(Elements(G),x->Order(x)=2);total:=0;bad:=0;markings:=[];
 for a in inv do for c in inv do if a<>c and a*c=c*a then for b in inv do
 rr:=[a,b,c];if Length(Set(rr))=3 and Subgroup(G,rr)=G then
 total:=total+1;Add(markings,rr);if not IP258(G,rr) then bad:=bad+1;fi;
 Need258(Order(a*b)>2 and Order(b*c)>2,"irreducible local pair");
 fi;od;fi;od;od;
 Need258(total=pair[2] and bad=0,"complete sggi to C local bridge");
 types:=[];aut:=AutomorphismGroup(G);left:=ShallowCopy(markings);
 while Length(left)>0 do rr:=left[1];orb:=Orbit(aut,rr,OnTuples);
 Add(types,[Order(rr[1]*rr[2]),Order(rr[2]*rr[3]),Order(Product(rr))]);
 left:=Filtered(left,x->not x in orb);od;
 Sort(types);Need258(types=pair[3],"all-sggi automorphism orbits");
 Print("SGGI_LOCAL ",Size(G)," TOTAL ",total," BAD_IP ",bad," TYPES ",types,"\n");
 od;
 # Exact counterexample to literal source Lemma3.3 maximality.
 G:=AlternatingGroup(5);z:=(1,2)(3,4);H:=Stabilizer(G,5);cs:=Centralizer(G,z);
 Need258(Size(cs)=4 and Size(H)=12 and IsSubgroup(H,cs) and H<>G,"centralizer maximality counterexample");
 Print("SOURCE_L33_COUNTEREXAMPLE PSL5 CENTRALIZER4 CONTAINED_A4_12\n");
 # Omitting IP in an arbitrary group really enlarges the local class.
 G:=DihedralGroup(IsPermGroup,12);inv:=Filtered(Elements(G),x->Order(x)=2);rr:=fail;
 for a in inv do for b in inv do for c in inv do
 if Length(Set([a,b,c]))=3 and a*c=c*a and Subgroup(G,[a,b,c])=G
 and Order(a*b)>2 and Order(b*c)>2 and not IP258(G,[a,b,c]) then rr:=[a,b,c];break;fi;
 od;if rr<>fail then break;fi;od;if rr<>fail then break;fi;od;
 Need258(rr<>fail,"no blanket sggi implies C assertion");
 Print("DIHEDRAL_NON_C_SGGI_CONTROL TYPE ",List([1,2],i->Order(rr[i]*rr[i+1])),"\n");
 # A centre-free group with the common-centralizer condition can have five
 # distinct involutory generators if the string relations are omitted.
 G:=AlternatingGroup(5);
 rr:=[(1,2)(3,4),(1,3)(2,4),(1,2)(3,5),(1,3)(2,5),(1,4)(2,5)];
 Need258(Length(Set(rr))=5 and Group(rr)=G,"no-string five-generator control");
 inv:=Filtered(Elements(G),x->Order(x)=2);
 for a in inv do for b in inv do if a*b<>b*a then
 Need258(Size(Centralizer(G,Subgroup(G,[a,b])))<=2,"pair-centralizer premise in control");
 fi;od;od;
 Need258(rr[1]*rr[3]<>rr[3]*rr[1],"actual omitted string relation");
 Print("RANK_BOUND_NO_STRING_FALSE_WEAKENING_A5_FIVE_MARKERS_REJECTED\n");
 Print("PHASE258_LOCAL_PASS\n");
end;
Run258();QUIT;
\end{CompleteProgram}
% END COMPLETE FILE phase258_local_sggi.g
\subsection*{Families and source corrections}
\noindent\texttt{phase258\_family\_check.g}
% BEGIN COMPLETE FILE phase258_family_check.g
\begin{CompleteProgram}
# Exact bounded projective-line models; uniform proofs are in the manuscript.
Need258:=function(b,s) if not b then Error(s);fi;end;
Projectivity258:=function(M,p)
 local out,x,den,y;
 out:=[];
 for x in [0..p] do
 if x=p then
 if M[2][1] mod p=0 then y:=p;else y:= (M[1][1]*PowerModInt(M[2][1],-1,p)) mod p;fi;
 else den:=(M[2][1]*x+M[2][2]) mod p;
 if den=0 then y:=p;else y:=((M[1][1]*x+M[1][2])*PowerModInt(den,-1,p)) mod p;fi;
 fi;Add(out,y+1);od;return PermList(out);
end;
String258:=function(rr)
 local i,j;
 if Length(Set(rr))<>Length(rr) or not ForAll(rr,x->Order(x)=2) then return false;fi;
 for i in [1..Length(rr)] do for j in [i+1..Length(rr)] do
 if j>i+1 and rr[i]*rr[j]<>rr[j]*rr[i] then return false;fi;
 od;od;return true;
end;
Run258:=function()
 local p,rr,G,orders,a,C,CC,src,changed,H,subs,inv,x,y,z,triple,D,cent,ext,kind,tgt,ms,dual,ds,i,j;
 for p in [13,37,61] do
 ms:=[[[0,1],[1,0]],[[1,0],[1,-1]],[[-1,0],[0,1]],[[1,-2],[2,-1]]];
 Need258(ForAll(ms,M->PowerModInt(DeterminantMat(M),(p-1)/2,p)=1),"constant determinant-square membership");
 rr:=List(ms,M->Projectivity258(M,p));
 G:=Group(rr);Need258(String258(rr),"constant string matrices");
 Need258(Size(G)=p*(p^2-1)/2 and IsSimpleGroup(G),"constant PSL image");
 Need258(Size(Subgroup(G,rr{[1..3]}))=Size(G) and Size(Subgroup(G,rr{[2..4]}))=Size(G),"two full ends");
 Need258(Order(rr[1]*rr[2])=3 and Order(rr[2]*rr[3])=p,"source3,p type");
 a:=First([1..p-1],x->x^2 mod p=p-1);C:=Projectivity258([[-a,0],[0,1]],p);
 src:=List([[[0,1],[-1,0]],[[a,0],[1,-a]],[[-a,0],[0,a]]],M->Projectivity258(M,p));
 # GAP right-action permutations reverse composition of the displayed column matrices.
 Need258(List(src,s->s^C)=rr{[1..3]},"literal source basis-change adapter");
 Print("CONSTANT_FAMILY ",p," ORDER ",Size(G)," TYPE ",List([1..3],i->Order(rr[i]*rr[i+1]))," SOURCE_CONJUGACY_PASS\n");
 od;
 # The smallest omitted S4 example is literal, avoiding a subgroup search.
 rr:=[(1,2)(3,24)(4,23)(5,22)(6,21)(7,20)(8,19)(9,18)(10,17)(11,16)(12,15)(13,14),
 (1,5)(2,7)(3,21)(4,11)(6,17)(8,23)(9,13)(10,24)(12,14)(15,18)(16,19)(20,22),
 (1,4)(2,23)(3,9)(5,13)(6,16)(7,10)(8,15)(11,21)(12,19)(14,22)(17,20)(18,24),
 (1,10)(2,17)(3,22)(4,18)(5,24)(6,7)(8,13)(9,23)(11,15)(12,16)(14,19)(20,21)];
 G:=PSL(2,23);Need258(ForAll(rr,x->x in G) and Group(rr)=G and String258(rr),"literal PSL23 marking");
 H:=Subgroup(G,rr{[1..3]});Need258(Size(H)=24 and IsomorphismGroups(H,SymmetricGroup(4))<>fail,"PSL23 S4 end");
 Need258(Subgroup(G,rr{[2..4]})=G and Index(G,H)=253,"omitted253facet count");
 Need258(List([1..3],i->Order(rr[i]*rr[i+1]))=[3,4,12],"literal omitted-family type");
 Print("OMITTED_S4 PSL23 ORDER6072 FACETS253 VERTICES1 TYPE3,4,12 PASS\n");
 # All triples of each chosen maximal subgroup are searched without an IP input.
 for ds in [[41,60,5,3],[47,24,3,4],[73,24,3,4]] do
 p:=ds[1];G:=PSL(2,p);
 subs:=MaximalSubgroupClassReps(G);
 if ds[2]=60 then tgt:=AlternatingGroup(5);else tgt:=SymmetricGroup(4);fi;
 H:=First(subs,x->Size(x)=ds[2] and IsomorphismGroups(x,tgt)<>fail);
 Need258(H<>fail,"retained maximal subgroup exists in diagnostic field");
 inv:=Filtered(Elements(H),x->Order(x)=2);triple:=fail;
 for x in inv do for y in inv do for z in inv do
 if Length(Set([x,y,z]))=3 and x*z=z*x and Order(x*y)=ds[3] and Order(y*z)=ds[4] and Subgroup(H,[x,y,z])=H then triple:=[x,y,z];break;fi;
 od;if triple<>fail then break;fi;od;if triple<>fail then break;fi;od;
 Need258(triple<>fail,"source local marking exists");D:=Subgroup(H,triple{[1,2]});cent:=Centralizer(G,D);
 Need258(Size(cent)=2 and Size(Centralizer(H,D))=1,"halfturn leaves facet");
 ext:=First(Elements(cent),x->Order(x)=2);rr:=Concatenation(triple,[ext]);
 Need258(String258(rr) and Subgroup(G,rr)=G and Subgroup(G,rr{[2..4]})=G,"maximal facet extension");
 Print("PROPER_END_FAMILY ",p," FACET ",Size(H)," FULL ",Size(G)," TYPE ",List([1..3],i->Order(rr[i]*rr[i+1]))," PASS\n");
 od;
 # A literal self-dual Class1 marking and its actual degree14 conjugator.
 rr:=[(3,9)(4,10)(5,11)(6,12)(7,13)(8,14),(1,4)(2,14)(3,9)(5,13)(6,11)(7,12),
 (1,2)(3,7)(4,6)(8,14)(9,13)(10,12),(1,2)(4,14)(5,13)(6,12)(7,11)(8,10)];
 G:=PSL(2,13);dual:=(1,3)(2,9)(4,7)(5,8)(6,12)(10,11)(13,14);
 Need258(ForAll(rr,x->x in G) and Group(rr)=G and String258(rr),"literal selfdual marking");
 Need258(Subgroup(G,rr{[1..3]})=G and Subgroup(G,rr{[2..4]})=G,"selfdual Class1 ends");
 Need258(List([1..3],i->Order(rr[i]*rr[i+1]))=[3,7,3] and List(rr,x->x^dual)=Reversed(rr),"literal selfdual reversal");
 Need258(List(rr,x->x^(dual*(1,2)))<>Reversed(rr),"altered conjugator rejected");
 Print("SELFDUAL_CLASS1 PSL13 ORDER1092 TYPE3,7,3 CONJUGATOR_PASS ALTERED_REJECTED\n");
 ext:=Concatenation(rr,[()]);
 Need258(Length(Set(ext))=5 and Group(ext)=G and ForAll(ext,x->x*x=()),"rankfive identity-marker countermodel");
 for i in [1..5] do for j in [1..5] do if j>i+1 then
 Need258(ext[i]*ext[j]=ext[j]*ext[i],"rankfive string commutations");fi;od;od;
 Need258(not String258(ext),"nonidentity condition is essential");
 Print("RANK_BOUND_ALLOW_IDENTITY_FALSE_WEAKENING_RANK5_REJECTED\n");
 Print("PHASE258_FAMILY_PASS\n");
end;
Run258();QUIT;
\end{CompleteProgram}
% END COMPLETE FILE phase258_family_check.g
\end{document}
